% !TeX program = xelatex
\documentclass{article}

% ============================================================
% Package selection
% ============================================================

\newif\ifUseXePersian
\UseXePersianfalse
% \UseXePersiantrue

% ============================================================
% Common packages
% ============================================================

\usepackage[a4paper,margin=25mm]{geometry}
\usepackage[hidelinks]{hyperref}

% ============================================================
% Persian package
% ============================================================

\ifUseXePersian
\usepackage{xepersian}
\newcommand{\ActivePackageName}{XePersian}
\else
\usepackage{unipersian}
\newcommand{\ActivePackageName}{UniPersian}
\fi

% ============================================================
% Fonts
% ============================================================

\settextfont{Amiri}
\setlatintextfont{TeX Gyre Termes}

% ============================================================
% Document
% ============================================================

\begin{document}

    \begin{center}
        {\Large آزمون دستور \lr{\textbackslash Latincite}}

        \medskip
        بستهٔ فعال: \lr{\ActivePackageName}
    \end{center}

    \hrule
    \bigskip

    \tableofcontents

    \listoffigures

    % ============================================================
    \section{Basic Citations}

    این یک ارجاع معمولی است:
    \cite{knuth}.

    این یک ارجاع لاتین تکی است:
    \Latincite{knuth}.

    این هم ارجاع به دو منبع است:
    \Latincite{knuth,lamport}.

    ترتیب معکوس دو منبع:
    \Latincite{lamport,knuth}.

    سه منبع با ترتیب متفاوت:
    \Latincite{mittelbach,knuth,lamport}.

    % ============================================================
    \section{Optional Notes}

    یادداشت فارسی ساده:
    \Latincite[صفحهٔ ۲۳]{knuth}.

    یادداشت فارسی با ویرگول:
    \Latincite[صفحهٔ ۲۳، بند ۴]{knuth}.

    یادداشت ترکیبی فارسی و انگلیسی:
    \Latincite[برای توضیح بیشتر در \lr{The TeXbook}، صفحهٔ ۲۳]{knuth}.

    یادداشت طولانی:
    \Latincite[برای مطالعهٔ بیشتر، به فصل دوم و صفحات ۲۳ تا ۳۱ مراجعه کنید]{knuth}.

    یادداشت لاتین:
    \Latincite[see page 23]{knuth}.

    % ============================================================
    \section{Multiple and Repeated Keys}

    کلیدها با فاصله:
    \Latincite{knuth, lamport}.

    کلیدهای تکراری:
    \Latincite{knuth,knuth}.

    سه کلید با فاصله:
    \Latincite{knuth, lamport, mittelbach}.

    تکرار در ترکیب با یادداشت:
    \Latincite[صفحات ۲۳ و ۲۴]{knuth,lamport}.

    % ============================================================
    \section{Interaction with Normal Citations}

    ابتدا ارجاع معمولی و سپس ارجاع لاتین:
    \cite{knuth}، \Latincite{knuth,lamport}.

    ابتدا ارجاع لاتین و سپس ارجاع معمولی:
    \Latincite{lamport,knuth}، \cite{mittelbach}.

    چند ارجاع متوالی:
    \Latincite{knuth}؛
    \Latincite[صفحهٔ ۴۰]{lamport}؛
    \cite{mittelbach}.

    % ============================================================
    \section{Citations in Moving Arguments
        \texorpdfstring{\Latincite{mittelbach}}{}}

    ارجاع در عنوان قسمت باید در فهرست مطالب نیز درست نمایش داده شود.

    \begin{figure}[htbp]
        \centering
        \rule{0.65\linewidth}{18mm}
        \caption{آزمون ارجاع در عنوان شکل
            \Latincite[صفحهٔ ۱۵]{lamport}}
        \label{fig:citation}
    \end{figure}

    ارجاع به شکل:
    \hyperref[fig:citation]{شکل بالا}.

    % ============================================================
    \section{Citations in Footnotes}

    ارجاع در متن پانویس%
    \footnote{
        این پانویس شامل یک ارجاع لاتین است:
        \Latincite[صفحهٔ ۲۳]{knuth}.
        متن فارسی پس از ارجاع نیز باید درست نمایش داده شود.
    }

    ارجاع دیگری در پانویس:
    \footnote{
        دو منبع:
        \Latincite{lamport,mittelbach}.
    }

    % ============================================================
    \section{Citations in Tables}

    \begin{center}
        \begin{tabular}{|c|c|c|}
            \hline
            منبع & ارجاع معمولی & ارجاع لاتین \\
            \hline
            Knuth & \cite{knuth} & \Latincite{knuth} \\
            \hline
            Lamport & \cite{lamport} & \Latincite{lamport} \\
            \hline
            Mittelbach & \cite{mittelbach} &
            \Latincite{mittelbach} \\
            \hline
        \end{tabular}
    \end{center}

    % ============================================================
    \section{Citations in Minipages}

    \begin{minipage}{0.8\linewidth}
        این یک محیط مستقل است.

        ارجاع تکی:
        \Latincite{knuth}.

        ارجاع چندتایی:
        \Latincite{lamport,mittelbach}.

        یادداشت فارسی:
        \Latincite[صفحهٔ ۵۰]{lamport}.
    \end{minipage}

    % ============================================================
    \section{Normal Text After Citations}

    \Latincite{knuth,lamport}.

    این متن فارسی باید پس از ارجاع همچنان
    با جهت و قلم صحیح ادامه پیدا کند.

    \Latincite[صفحهٔ ۲۳]{knuth}.

    این هم متن فارسی پس از ارجاع دارای یادداشت.

    % ============================================================
    \begin{thebibliography}{99}

        \bibitem{knuth}
        \lr{%
            Donald E. Knuth.
            \emph{The TeXbook}.
            Addison-Wesley, 1984.
        }

        \bibitem{lamport}
        \lr{%
            Leslie Lamport.
            \emph{LaTeX: A Document Preparation System}.
            Addison-Wesley, 1994.
        }

        \bibitem{mittelbach}
        \lr{%
            Frank Mittelbach et al.
            \emph{The LaTeX Companion}.
            Addison-Wesley, 2004.
        }

    \end{thebibliography}

\end{document}